% !TEX root = ../main.tex

\chapter{绪论}

\section{研究背景与意义}

\subsection{水下机器人概述}
水下机器人是一种能够在水下环境中执行各种任务的自主或遥控设备。它们被广泛应用于海洋科学研究、资源勘探、环境监测、军事侦察等领域。
随着技术的不断进步，水下机器人的性能和功能也在不断提升，使其在复杂水域中的应用变得更加广泛和高效。 
在“海洋强国”战略部署与海洋经济高质量发展的双重驱动下，浅水水域的水产养殖智能化监测、近海风电运维及水下工程维护已成为深蓝经济增长的关键领域。
水下机器人已经成为实现这些任务的重要工具，能够在复杂的水下环境中执行各种任务，如数据采集、环境监测、资源勘探等。

根据水下机器人的通信方式、供能方式、任务类型等不同，可以将水下机器人分为多种类型，如自主水下机器人（AUV）、遥控水下机器人（ROV）等。
每种类型的水下机器人都有其独特的设计和功能，以适应不同的应用需求。
前者是无线缆的水下机器人，主要通过预先编写的程序实现在水下的自主导航与操作\cite{Kang2023}，
后者则是通过有线连接与操作人员进行通信和控制，适用于需要实时监控和操作的任务。两者的详细对比参考表\ref{tab:rov-auv-compare}。

\begin{table}[htbp]
    \centering
    \caption{ROV与AUV对比}
    \label{tab:rov-auv-compare}
    \begin{tabular}{cccc}
        \toprule
        对比项 & ROV & AUV \\
        \midrule
        控制方式 & 通过线缆实时遥控 & 自主导航与决策 \\
        供电方式 & 线缆供电 & 电池供电 \\
        应用场景 & 精细操作任务 & 大范围探测任务 \\
        作业范围 & 受脐带缆长度限制 & 受电池续航限制 \\
        \bottomrule
    \end{tabular}
\end{table}

\subsection{浅水水域作业现状与挑战}

\subsection{水下机器人控制系统概念}

\subsection{研究意义}

\section{国内外研究现状}

\subsection{无模型控制}

\subsection{模型控制}

\subsection{基于学习的控制}

\section{本文主要研究内容与目标}

\section{论文结构安排}

\section{本章小结}
